\begin{lemma}
\label{lemma-total-ring-fractions}
Let $R$ be a ring.
Let $S \subset R$ be a multiplicative subset consisting of nonzerodivisors.
Then $Q(R) \cong Q(S^{-1}R)$.
In particular $Q(R) \cong Q(Q(R))$.
\end{lemma}

\begin{proof}
Let $T\subset R$ be the set of nonzerodivisors.
Since $S\subset T$, the map $R\to S^{-1}R$ is injective.
An element $x=r/f$ of $S^{-1}R$, with $r\in R$ and $f\in S$,
is a nonzerodivisor if and only if $r\in T$.
Indeed, if $x$ is a nonzerodivisor and $ra=0$ in $R$, then
$x(a/1)=0$, hence $a/1=0$, and injectivity gives $a=0$.
Conversely, if $r\in T$ and $x(a/s)=0$, then
$ura=0$ for some $u\in S$. Both $u$ and $r$ are nonzerodivisors,
so $a=0$ and $a/s=0$.
The two ring maps are
$$
Q(R)\longrightarrow Q(S^{-1}R),\qquad
\frac{r}{t}\longmapsto\frac{r/1}{t/1},
$$
and
$$
Q(S^{-1}R)\longrightarrow Q(R),\qquad
\frac{r/s}{a/u}\longmapsto\frac{ru}{sa},
\quad s,u\in S,\quad a\in T.
$$
The first is induced by $R\to S^{-1}R$ because each $t\in T$
remains a nonzerodivisor. The second is induced by $S^{-1}R\to Q(R)$:
it sends every nonzerodivisor $a/u$ to a unit with inverse $u/a$.
Thus both formulas are well-defined by the localization universal
properties. Their composites fix each element of $R$ and every
adjoined inverse; explicitly, $r/t$ returns to $r/t$, and
$(r/s)/(a/u)$ returns to $(ru/1)/(sa/1)$, the original fraction.
Taking $S=T$ gives the final assertion, with the same maps.
\end{proof}